\chapter{Introduction} \label{Intoduction}
This thesis looks at the development of a software system for the Danish Environmental Protection Agency, with the purpose of helping them to uphold the regulations put in place regarding pre-harvest crop desiccant use. It will cover the necessary agricultural knowledge regarding Danish potato production and crop desiccation methods, as well as exploring how satellite imagery can be used to monitor crop health over time, and by extension, whether crop desiccant usage can be monitored in relation to the desiccation time of crops. The remaining sections in the introduction will provide further insight into the case presented by the Danish Environmental Protection Agency, the problem the authors are attempting to solve, how the case was approached and important terminology. A reading guide is presented at the end to help navigate the thesis report.

\section{Collaboration with the Danish Environmental Protection Agency}
The \acrfull{DEPA} (da: Miljøstyrelsen)(DEPA) is a government agency that deals with a variety of tasks related to Danish business, industry, agriculture, and consulting companies and business organizations, to further Danish interests while ensuring environmental progress\cite{MiljøstyrelsenOpgaver}. One such task is upholding the regulations put in place in regards to pre-harvest crop desiccant usage in the agricultural sector, to prevent pollution of nearby ecosystems and the ground water. The \acrshort{DEPA} contacted The University of Southern Denmark (SDU) through a representative with a project case, that could help uphold the pre-harvest crop desiccant regulations. Thus, this representative was the primary contact person from the DEPA for the authors throughout the project.


\section{Context}
The following section is an introduction into the case presented by the DEPA representative and gives further context into problem they are facing and how it could potentially be solved using satellite imagery. Additionally, the case presentation can be found in appendix(INDSÆTPP).

Pre-harvest crop desiccants are commonly used for wilting crops on fields. Ensuring that all plants are dead and wilted will significantly diminish the bacteria and parasites that are left to carry over to the next year, thus improving the fundamental conditions for future harvests\cite{europaAgriculturePesticidesplant}. But pesticides have been shown to seep into the ground water, which poses significant health implications\cite{pesticiderGrundvand} and to be damaging to nearby wildlife ecosystems\cite{Middeldatabasen}. In order to restrict the rate of this pollution, regulations made by the DEPA have been put in place for how much and how often pesticides are allowed to be sprayed. To uphold these regulations, mandated annual inspections are performed. But considering that 60.4\% of Denmark's landscape is utilised for the agricultural sector\cite{Lindsgaard2017-em}, that the number of mandated annual inspections are below 1000 a year (746 inspections in 2019)\cite{MiljoeOgFoedevareministeriet}, and that the inspections are random rather than being informed by some observations, the chances of getting caught for an overuse of pesticides are slim. 

A map of the NDVI index can be produced from data collected by the Sentinel 1 through 6 satellites from the Copernicus Expansion Missions\cite{noauthor_undated-bb}. Normalized Difference Vegetation Index (NDVI) is a common index of vegetation health used in remote sensing\cite{GISGeography2017-cr}. Healthy vegetation reflects more near-infrared light and absorbs more visible red light, resulting in a higher NDVI value. In contrast, stressed vegetation absorbs more near-infrared light and reflects more red light, leading to a lower NDVI value. The scale of NDVI goes from -1 to 1, where a 1 would be a healthy plant and 0 would indicate no healthy vegetation or an urbanised area. A negative NDVI index would indicate an ocean\cite{GISGeography2017-cr}. With this information one should be able to deduce whether there are stressed or wilting crops on fields and how quickly the wilting happens, if monitored over time. 

A quicker than usual wilting process may be an indication of an overuse of pesticides, which is in violation of the regulations, and such an observation might warrant an inspection. By inspecting those crops that are more likely to have been wilted using pesticides, the hit rate of inspections should rise. The number of farmers that will get fined for this overuse will produce additional funding that can lead to a more balanced risk scenario for the farmers, resulting in less pesticide use and thus less pollution of the ground water and nearby wildlife ecosystems.


\section{Problem}
What would it take to implement a system for finding inspection site candidates, and how efficient would it be in finding overuse of pre-harvest crop desiccants?
\vspace{1em}

\noindent In particular,
\begin{enumerate}
    \item \textbf{Model:} How can satellite data be analysed to define a statistical model of crop wilting where NDVI is defined as a function over time? Such a model should be parameterised over geography and type of crop.
    \item \textbf{Validation:} Can such a model be used to detect crop desiccation rates that fall outside the model and detect pre-harvest crop desiccant use?
    \item \textbf{Scale:} How can such a model be implemented efficiently, for all crop fields in Denmark?
\end{enumerate}

\section{Approach}
Historical observations of normal crop desiccation timelines will be used to train and build a models of a single crop type commonly produced in Denmark.
One model will be produced for normal wilting and another for desiccant accelerated wilting. Open-access Sentinel data will be used for this, and ground truth from DEPA will be used to differentiate between normal and accelerated wilting.
This will allow the observed wilting process of any field to be compared to the normal model and the accelerated model, and based of such comparisons a likelihood of crop desiccant use can be determined.
Cross-validation using the ground truth from DEPA will be performed to evaluate to which degree the model-pair is capable of differentiating between normal and accelerated wilting. Once such differentiation is possible, finding outliers from both models, which could also be an indication of pesticide overuse, is possible.

Starting off, the model will be trained to detect accelerated wilting in potato plants, by request of the DEPA. Then, to scale the model, other crop types will be chosen for training and validation. The whole model will then be deployed on a cloud service with big data techniques to ensure that the DEPA can scale and launch the monitoring service for different crop types and find outlier from the model, prompting a site inspection.


\section{Terminology}
Vi udpensler allerede NDVI ret godt i context afsnittet, men vil gerne lige snakke om hvilke andre terminologier vi evt skal have fat i. Evt sådan noget "hvad mener vi, når vi siger "crop wilting" eller "model for crop wilting over time" eller whatever"

\section{Reading Guide}
Skrives senere